% GENERATED FILE. Edit canonical lessons, then run npm run generate:books.
% canonical-source-hash: fnv1a64:8300c35259b0e15f
% canonical-lessons: GU-C383-prayatna, GU-C383-kosis, GU-C383-utsah, GU-C383-madvu, GU-C383-jhujhvu

\chapter{Effort}
\label{ch:gu-effort}
\hlchaptermodality{383}

\begin{chapteropening}
\textbf{By the end of this chapter:} \emph{I can say what effort it will take to succeed.}

Complete the last lesson of chapter 383: I can say what effort it will take to succeed.
\end{chapteropening}

\section[\texorpdfstring{\gu{પ્રયત્ન} (prayatna)}{પ્રયત્ન (prayatna)}]{\gu{પ્રયત્ન} (prayatna) — an attempt}

\label{lesson:GU-C383-prayatna}



\begin{quote}
Before the new one: say the Gujarati for planning, then the Gujarati for to strive.
\end{quote}

\subsection*{You'll want to know: \gu{પ્રયત્ન}}
\textbf{\gu{પ્રયત્ન}} — \emph{prayatna} — \textquotedblleft{}an attempt\textquotedblright{}.

\begin{quote}
\gu{જો} \gu{પહેલા} \gu{પ્રયત્નમાં} \gu{તમે} \gu{હારશો} \gu{તો} \gu{ફરી} \gu{પ્રયત્ન} \gu{કરો}.
\end{quote}

\begin{quote}
\emph{jo pahelā prayatnamā̃ tame hārśo to pharī prayatna karo.}
\end{quote}

\begin{quote}
If you lose at the first attempt, try again.
\end{quote}

A masculine noun.

\begin{cousinweb}
\textbf{\gu{પ્રયત્ન}} (\emph{prayatna}) is a tatsama, borrowed from Sanskrit \emph{prayatna}, \textquotedblleft{}persevering effort, exertion\textquotedblright{} (Monier-Williams).
\end{cousinweb}

\begin{grammarlens}[title={What you've built}]
One more word for saying what happens if.
\end{grammarlens}

\subsection*{Your turn}
\begin{itemize}
\raggedright
  \item \emph{Say it:} \emph{prayatna}
  \item \emph{Say it:} \emph{prayatna}, once more
  \item \emph{Say it:} \emph{jo pahelā prayatnamā̃ tame hārśo to pharī prayatna karo.} — If you lose at the first attempt, try again.
  \item \emph{Recall:} say the Gujarati for planning, then the Gujarati for to strive, then say \emph{prayatna} again
\end{itemize}

\begin{tcolorbox}[breakable,colback=teal!4,colframe=teal!35!black,title={Before you move on}]
What does \gu{પ્રયત્ન} mean? (\textquotedblleft{}An attempt\textquotedblright{}.) Say it once more.
\end{tcolorbox}
\section[\texorpdfstring{\gu{કોશિશ} (kośiś)}{કોશિશ (kośiś)}]{\gu{કોશિશ} (kośiś) — an effort}

\label{lesson:GU-C383-kosis}



\begin{quote}
Before the new one: say the Gujarati for to strive, then the Gujarati for an attempt.
\end{quote}

\subsection*{You'll want to know: \gu{કોશિશ}}
\textbf{\gu{કોશિશ}} — \emph{kośiś} — \textquotedblleft{}an effort\textquotedblright{}.

\begin{quote}
\gu{જો} \gu{તમે} \gu{કોશિશ} \gu{કરશો} \gu{તો} \gu{હું} \gu{પણ} \gu{મદદ} \gu{કરીશ}.
\end{quote}

\begin{quote}
\emph{jo tame kośiś karśo to hũ paṇ madad karīś.}
\end{quote}

\begin{quote}
If you make the effort, I will help too.
\end{quote}

A feminine noun.

\begin{cousinweb}
\textbf{\gu{કોશિશ}} (\emph{kośiś}) is usually traced to Persian \emph{kōshish}, \textquotedblleft{}effort, striving\textquotedblright{}.
\end{cousinweb}

\begin{grammarlens}[title={What you've built}]
One more word for saying what happens if.
\end{grammarlens}

\subsection*{Your turn}
\begin{itemize}
\raggedright
  \item \emph{Say it:} \emph{kośiś}
  \item \emph{Say it:} \emph{kośiś}, once more
  \item \emph{Say it:} \emph{jo tame kośiś karśo to hũ paṇ madad karīś.} — If you make the effort, I will help too.
  \item \emph{Recall:} say the Gujarati for to strive, then the Gujarati for an attempt, then say \emph{kośiś} again
\end{itemize}

\begin{tcolorbox}[breakable,colback=teal!4,colframe=teal!35!black,title={Before you move on}]
What does \gu{કોશિશ} mean? (\textquotedblleft{}An effort\textquotedblright{}.) Say it once more.
\end{tcolorbox}
\section[\texorpdfstring{\gu{ઉત્સાહ} (utsāh)}{ઉત્સાહ (utsāh)}]{\gu{ઉત્સાહ} (utsāh) — enthusiasm}

\label{lesson:GU-C383-utsah}



\begin{quote}
Before the new one: say the Gujarati for an attempt, then the Gujarati for an effort.
\end{quote}

\subsection*{You'll want to know: \gu{ઉત્સાહ}}
\textbf{\gu{ઉત્સાહ}} — \emph{utsāh} — \textquotedblleft{}enthusiasm\textquotedblright{}.

\begin{quote}
\gu{જો} \gu{ખેલાડીઓમાં} \gu{ઉત્સાહ} \gu{હશે} \gu{તો} \gu{તેઓ} \gu{સારું} \gu{રમશે}.
\end{quote}

\begin{quote}
\emph{jo khelāḍīomā̃ utsāh haśe to teo sārũ ramśe.}
\end{quote}

\begin{quote}
If the players are keen, they will play well.
\end{quote}

A masculine noun.

\begin{cousinweb}
\textbf{\gu{ઉત્સાહ}} (\emph{utsāh}) is a tatsama, borrowed from Sanskrit \emph{utsāha}, \textquotedblleft{}power, effort, zeal\textquotedblright{} (Monier-Williams).
\end{cousinweb}

\begin{grammarlens}[title={What you've built}]
One more word for saying what happens if.
\end{grammarlens}

\subsection*{Your turn}
\begin{itemize}
\raggedright
  \item \emph{Say it:} \emph{utsāh}
  \item \emph{Say it:} \emph{utsāh}, once more
  \item \emph{Say it:} \emph{jo khelāḍīomā̃ utsāh haśe to teo sārũ ramśe.} — If the players are keen, they will play well.
  \item \emph{Recall:} say the Gujarati for an attempt, then the Gujarati for an effort, then say \emph{utsāh} again
\end{itemize}

\begin{tcolorbox}[breakable,colback=teal!4,colframe=teal!35!black,title={Before you move on}]
What does \gu{ઉત્સાહ} mean? (\textquotedblleft{}Enthusiasm\textquotedblright{}.) Say it once more.
\end{tcolorbox}
\section[\texorpdfstring{\gu{માંડવું} (mā̃ḍvũ)}{માંડવું (mā̃ḍvũ)}]{\gu{માંડવું} (mā̃ḍvũ) — to set about}

\label{lesson:GU-C383-madvu}



\begin{quote}
Before the new one: say the Gujarati for an effort, then the Gujarati for enthusiasm.
\end{quote}

\subsection*{You'll want to know: \gu{માંડવું}}
\textbf{\gu{માંડવું}} — \emph{mā̃ḍvũ} — \textquotedblleft{}to set about\textquotedblright{}.

\begin{quote}
\gu{જો} \gu{આપણે} \gu{આજથી} \gu{કામ} \gu{માંડીશું} \gu{તો} \gu{મહિનામાં} \gu{ઘર} \gu{તૈયાર} \gu{થશે}.
\end{quote}

\begin{quote}
\emph{jo āpaṇe ājthī kām mā̃ḍīśũ to mahināmā̃ ghar taiyār thaśe.}
\end{quote}

\begin{quote}
If we set to work today, the house will be ready within a month.
\end{quote}

Past, with \gu{તેણે} (he or she …): \gu{તેણે} \gu{નવો} \gu{ધંધો} \gu{માંડ્યો}, \emph{teṇe navo dhandho mā̃ḍyo}.

\begin{cousinweb}
\textbf{\gu{માંડવું}} (\emph{mā̃ḍvũ}) is inherited from Sanskrit \emph{maṇḍáyati}, \textquotedblleft{}adorns, decorates\textquotedblright{}; CDIAL lists the Gujarati verb, \textquotedblleft{}to arrange, to begin\textquotedblright{} (CDIAL 9741).
\end{cousinweb}

\begin{grammarlens}[title={What you've built}]
One more word for saying what happens if.
\end{grammarlens}

\subsection*{Your turn}
\begin{itemize}
\raggedright
  \item \emph{Say it:} \emph{mā̃ḍvũ}
  \item \emph{Say it:} \emph{mā̃ḍvũ}, once more
  \item \emph{Say it:} \emph{jo āpaṇe ājthī kām mā̃ḍīśũ to mahināmā̃ ghar taiyār thaśe.} — If we set to work today, the house will be ready within a month.
  \item \emph{Recall:} say the Gujarati for an effort, then the Gujarati for enthusiasm, then say \emph{mā̃ḍvũ} again
\end{itemize}

\begin{tcolorbox}[breakable,colback=teal!4,colframe=teal!35!black,title={Before you move on}]
What does \gu{માંડવું} mean? (\textquotedblleft{}To set about\textquotedblright{}.) Say it once more.
\end{tcolorbox}
\section[\texorpdfstring{\gu{ઝૂઝવું} (jhūjhvũ)}{ઝૂઝવું (jhūjhvũ)}]{\gu{ઝૂઝવું} (jhūjhvũ) — to battle on}

\label{lesson:GU-C383-jhujhvu}



\begin{quote}
Before the new one: say the Gujarati for enthusiasm, then the Gujarati for to set about.
\end{quote}

\subsection*{You'll want to know: \gu{ઝૂઝવું}}
\textbf{\gu{ઝૂઝવું}} — \emph{jhūjhvũ} — \textquotedblleft{}to battle on\textquotedblright{}.

\begin{quote}
\gu{જો} \gu{તે} \gu{રોગ} \gu{સામે} \gu{ઝૂઝશે} \gu{તો} \gu{સાજો} \gu{થશે}.
\end{quote}

\begin{quote}
\emph{jo te rog sāme jhūjhśe to sājo thaśe.}
\end{quote}

\begin{quote}
If he fights the illness, he will get better.
\end{quote}

Past (he …): \gu{તે} \gu{છેલ્લે} \gu{સુધી} \gu{ઝૂઝ્યો}, \emph{te chelle sudhī jhūjhyo}.

\begin{cousinweb}
\textbf{\gu{ઝૂઝવું}} (\emph{jhūjhvũ}) is inherited from Sanskrit \emph{yúdhyatē}, \textquotedblleft{}fights against\textquotedblright{}; CDIAL lists the Gujarati verb, \textquotedblleft{}to fight\textquotedblright{} (CDIAL 10502).
\end{cousinweb}

\begin{grammarlens}[title={What you've built}]
One more word for saying what happens if.
\end{grammarlens}

\subsection*{Your turn}
\begin{itemize}
\raggedright
  \item \emph{Say it:} \emph{jhūjhvũ}
  \item \emph{Say it:} \emph{jhūjhvũ}, once more
  \item \emph{Say it:} \emph{jo te rog sāme jhūjhśe to sājo thaśe.} — If he fights the illness, he will get better.
  \item \emph{Recall:} say the Gujarati for enthusiasm, then the Gujarati for to set about, then say \emph{jhūjhvũ} again
\end{itemize}

\begin{tcolorbox}[breakable,colback=teal!4,colframe=teal!35!black,title={Before you move on}]
What does \gu{ઝૂઝવું} mean? (\textquotedblleft{}To battle on\textquotedblright{}.) Say it once more.
\end{tcolorbox}
